\section{The consistency deduction}\label{sec:metatheory}

Theorem~\ref{thm:construction-summary} has two uses. For ordinary
relative consistency, we must obtain a choice ground from
$\Con(\ZF)$ and interpret the construction even when that ground
is externally ill founded. For the separate transitive-model result,
we retain the given transitive choice ground and evaluate names in
the usual way. We begin with the precise forcing assertion that
both deductions use.

\subsection{The internal output of the construction}

Write $\mathcal S=(\mathbb Q,\Gamma,\mathcal F)$ for a set
symmetric system: $\mathbb Q$ is a nonempty set poset, $\Gamma$
is a set group of automorphisms of $\mathbb Q$, and $\mathcal F$
is a normal filter of subgroups of $\Gamma$. Let
$\HS_{\mathcal S}$ be its class of hereditarily symmetric names.
We use a largest condition $1$. If necessary, adjoin a fresh
largest condition fixed by every automorphism; the original poset
remains dense.

The preceding construction is uniform over choice grounds. In
particular, Theorem~\ref{thm:construction-summary} gives the
following internal assertion in $\ZFC+V=L$: there is a nonempty
set symmetric system $\mathcal S$ such that
\begin{equation}\label{eq:construction-internal-forcing}
 1\forces_{\HS_{\mathcal S}}
          \PP\ \wedge\ \neg\AC\ \wedge\ \ACwo.
\end{equation}
Here the universal quantifiers of the three sentences range over
all internal HS names, with no name-rank cutoff. The standard
symmetric-model theorem supplies full $\ZF$ for this interpretation;
its formula-by-formula justification is given below.

The internal character of this particular construction matters.
The monoid is built by recursion along a set ordinal, and its
countable-base rearrangement uses an elementary submodel of a set
structure $H_\Theta$. Its normal forms, diagrams, permutations and
supports are sets. In the quotient argument,
Equation~\eqref{eq:normalized-local-names} bounds the local name
collections by ground sets. The assignment predicate
$\Phi$ in Equation~\eqref{eq:assignment-predicate} uses those
sets, fixed ground choices, and the ordinary forcing relation for
the ground poset $\mathbb P$ in an $\mathbb R$-extension.
Equation~\eqref{eq:quotient-bijection-name} then forms an actual
ground HS graph name. Thus the generic notation in that argument
expresses fixed-formula forcing assertions; it is not a construction
available only over externally transitive models.

Propositions~\ref{prop:presentation} and~\ref{prop:ordinal-choice}
supply the all-rank presentation and ordinal-indexed choice.
Together with the quotient graph, the internal $\ZF$ reduction in
Lemma~\ref{lemma:pp-reduction} gives $\PP$.
Proposition~\ref{prop:not-choice} supplies $\neg\AC$.
The quotient witnesses are produced below a strengthening of every
condition asserting the relevant hypothesis. This is the local-witness
form of existential forcing, made explicit in
Equation~\eqref{eq:symmetric-existential-density} below; no single
symmetric witness is required to work below every condition.

It remains to transfer this forcing assertion to a model obtained
from ordinary consistency. We first obtain a choice ground by the
constructible interpretation, then define the symmetric extension
without external recursion on an ill-founded membership relation.

\subsection{The constructible choice ground}

\begin{lemma}\label{lem:constructible-consistency}
For every axiom $\alpha$ of $\ZFC+V=L$, its relativization
$\alpha^L$ is provable in $\ZF$. Consequently
\[
 \Con(\ZF)\ \Longrightarrow\ \Con(\ZFC+V=L).
\]
Moreover, $\ZFC+V=L$ proves the generalized continuum hypothesis.
\end{lemma}

\begin{proof}
The constructible-universe theorem proves, in $\ZF$, the
relativization to $L$ of every $\ZF$ axiom, of Choice, and of
$V=L$. These are assertions for each individual axiom instance,
including the instances of Separation and Replacement. See
\cite[Theorem~4.7, Lemma~4.10, Corollary~4.13, and
Theorems~4.14 and~4.18]{Welch2020}.

Relativization commutes with Boolean connectives and preserves
first-order inference when quantifiers are restricted to the
nonempty interpreted domain. A finite contradiction proof from
$\ZFC+V=L$ would therefore yield, by relativizing each line and
using the finitely many required interpreted axioms, a
contradiction proof from $\ZF$. This proves the consistency
implication. The same constructible-universe theorem proves GCH
in $L$, hence proves GCH from $\ZFC+V=L$.
\end{proof}

In this choice ground the parameters of the construction are
\[
 \kappa=\aleph_1,\qquad \mu=\aleph_2,\qquad I=\aleph_3.
\]
These are a convenient specialization of the arbitrary-ground
parameters in Section~\ref{sec:forcing}. GCH gives
\[
 \mu^{\aleph_0}=(2^\kappa)^{\aleph_0}=\mu.
\]
Since $\kappa=\aleph_1$, this is the bound
$\mu^{<\kappa}=\mu$ used in the construction. CH also gives
$\kappa^{<\kappa}=\kappa$ here, but that additional equality is
not an assumption of the arbitrary-ground argument. The required
general bounds were proved in
Equation~\eqref{eq:ground-countable-arithmetic} and its preceding
paragraph. The set constructions and ground wellorderings used
earlier are available in $\ZFC+V=L$. In particular, $\Theta$ in
the monoid rearrangement need only be a sufficiently large regular
cardinal; no inaccessible cardinal or satisfaction predicate for
the universe is required.

\subsection{Symmetric forcing as a first-order interpretation}

We now justify the interpretation machinery used with
Equation~\eqref{eq:construction-internal-forcing}. The point is
not only to obtain the elementary set axioms: every instance of
Separation and Replacement must hold for arbitrary-rank names.

\begin{lemma}[The symmetric forcing schema]
\label{lem:symmetric-forcing-schema}
For each fixed formula $\varphi$ in the language of set theory,
the relation
\[
 p\forces_{\HS_{\mathcal S}}\varphi(\tau_1,\ldots,\tau_n)
\]
is definable in $\ZFC$, with $\mathcal S$ as a parameter and
$\tau_1,\ldots,\tau_n\in\HS_{\mathcal S}$.
It satisfies the forcing rules for first-order logic with equality
and the symmetry lemma. For each axiom $\alpha$ of $\ZF$,
including each Separation and Replacement instance, $\ZFC$
proves
\[
 \mathcal S\text{ is a set symmetric system}
 \quad\Longrightarrow\quad
 1\forces_{\HS_{\mathcal S}}\alpha.
\]
\end{lemma}

\begin{proof}
Atomic forcing is defined by the usual recursion on names.
For a fixed formula, extend these atomic definitions through its
finitely many connectives and quantifiers, restricting quantified
names to $\HS_{\mathcal S}$. For example, define
\[
 D_{\varphi,\vec\tau}
 =\{q\in\mathbb Q:(\exists\sigma\in\HS_{\mathcal S})\
              q\forces_{\HS_{\mathcal S}}\varphi(\sigma,\vec\tau)\}.
\]
Existential forcing is characterized by
\begin{equation}\label{eq:symmetric-existential-density}
 p\forces_{\HS_{\mathcal S}}\exists x\,\varphi(x,\vec\tau)
 \quad\Longleftrightarrow\quad
 D_{\varphi,\vec\tau}\text{ is dense below }p.
\end{equation}
The set $D_{\varphi,\vec\tau}$ exists by Separation from $\mathbb Q$.
The class of names and the hereditary symmetry predicate are
definable by set-theoretic recursion. Thus all the displayed
relations are ordinary first-order predicates. This construction
is a schema indexed externally by formulas; it asserts no
single satisfaction predicate for all formulas.

The logical and symmetry rules follow by induction on a fixed
formula, with the internal name recursion supplying the atomic
case. The resulting symmetric-model theorem and its forcing
relation are standard; see
\cite[Equations~(15.39)--(15.41) and Lemma~15.51]{Jech2003}
and \cite[Section~2.1, p.~4]{Karagila2026}.
We recall the parts that ensure the full axiom schemata.

Fix an arbitrary formula $\varphi(x,\vec a)$ and
hereditarily symmetric names $\sigma,\vec\tau$ for its set and
parameters. A Separation name uses the pairs
$(\rho,q)\in\dom(\sigma)\times\mathbb Q$ for which
\[
 q\forces_{\HS_{\mathcal S}}
       \rho\in\sigma\ \wedge\ \varphi(\rho,\vec\tau).
\]
The intersection of the stabilizers of $\sigma$ and the finitely
many parameter names fixes this name. That intersection belongs
to $\mathcal F$, and every member name is hereditarily symmetric.
The atomic forcing rules and
Equation~\eqref{eq:symmetric-existential-density} verify that it
is the required subset. No restriction on $\varphi$ or on the
rank of $\sigma$ has been imposed.

Write $\operatorname{nrk}(\eta)$ for the name-rank
\[
 \operatorname{nrk}(\eta)=
 \sup\{\operatorname{nrk}(\zeta)+1:(\zeta,r)\in\eta\}.
\]
For Replacement, fix a formula $\psi(x,y,\vec a)$, a domain
name $\sigma$, and parameter names $\vec\tau$. For each pair
$(\rho,q)\in\dom(\sigma)\times\mathbb Q$ for which some
hereditarily symmetric name $\eta$ satisfies
\[
 q\forces_{\HS_{\mathcal S}}\psi(\rho,\eta,\vec\tau),
\]
take the least possible name-rank of such an $\eta$.
This defines an ordinal-valued function on a ground set.
Replacement in the ground bounds all these ranks strictly below
one ordinal $\beta$. All $\mathbb Q$-names of name-rank less
than $\beta$ form a set. Consequently
\[
 \dot B_\beta=
 \{(\eta,1):\eta\in\HS_{\mathcal S},\
                         \operatorname{nrk}(\eta)<\beta\}
\]
is a set name. Every automorphism preserves name-rank and
hereditary symmetry, so $\dot B_\beta$ is fixed by the whole
group and belongs to $\HS_{\mathcal S}$.

Below a condition forcing totality of the relation on $\sigma$,
local witnesses from
Equation~\eqref{eq:symmetric-existential-density} can therefore
be replaced, at the same witnessing conditions, by names occurring
in $\dot B_\beta$. This gives a set bounding the required
witnesses. When the relation is functional, Separation in that
set gives its image. The bound $\beta$ depends on the particular
domain name and formula; it is not a bound on the universe of
names.

For Power Set, normalize a name for a subset of $\sigma$ using
member names from $\dom(\sigma)$ and the forcing tests for
membership in that subset. The normalized names lie in the
ground set $\mathcal P(\dom(\sigma)\times\mathbb Q)$.
Use all hereditarily symmetric candidates, under the conditions
forcing that they are subsets of $\sigma$. The stabilizer of
$\sigma$ permutes this collection, giving a hereditarily
symmetric power-set name. Every symmetric subset has such a
normalization, so the name supplies the entire internal power set.
Check names and the canonical pairing and union constructions
supply the remaining elementary existence axioms.

The usual internal name-rank argument verifies Foundation.
Equivalently, the interpreted hereditary class is membership-closed
inside the ordinary forcing extension: an element of an
HS-valued set has an active member name of an HS name.
Extensionality and Foundation then inherit their witnesses from
the ordinary extension. This argument concerns first-order
Foundation and does not assume that an arbitrary model is
externally well founded.

All these name constructions and forcing verifications are proofs
in the ground theory, for the chosen formula. They prove the
asserted schema. They use local existential witnesses rather
than a maximum principle for symmetric forcing; that maximum
principle need not hold
\cite[Proposition~10.26]{KaragilaNotes2026}.
\end{proof}

\subsection{The extension of an arbitrary countable model}

Let $M=(|M|,E^M)$ be a countable, possibly ill-founded model of
$\ZFC$. Suppose $M$ contains a set symmetric system
$\mathcal S$. Its internal poset determines the external set
\[
 Q^M=\{p\in|M|:M\models p\in\mathbb Q\},
\]
with the order interpreted by $M$. An $M$-generic filter means
a filter on this poset meeting the interpretation of every
$D\in|M|$ for which $M$ asserts that $D$ is dense in
$\mathbb Q$.

Such a filter exists through any prescribed condition. Enumerate
these dense-set codes externally and choose a descending sequence
meeting them successively. Internal density really supplies an
external strengthening in $Q^M$, since its witnesses are elements
of the domain of $M$. The upward closure of the chosen sequence
is the required generic filter $G$. No closure property of
$\mathbb Q$ and no well-foundedness of $E^M$ are needed here.

Let $\mathcal N^M$ be the external set of elements of $M$ which
$M$ regards as $\mathbb Q$-names. Define
\begin{align}
 \sigma\sim_G\tau
 &\quad\Longleftrightarrow\quad
   (\exists p\in G)\ M\models p\forces\sigma=\tau,
 \label{eq:internal-name-equivalence}\\
 [\sigma]_G\ E_G\ [\tau]_G
 &\quad\Longleftrightarrow\quad
   (\exists p\in G)\ M\models p\forces\sigma\in\tau.
 \label{eq:internal-name-membership}
\end{align}
The equality and substitution laws proved internally in $M$,
together with directedness of $G$, make $\sim_G$ an equivalence
relation and make $E_G$ independent of representatives.
The quotient of all internal names is the ordinary extension
structure $M[G]$. Its substructure $W$ consists of the equivalence
classes having a representative which $M$ regards as hereditarily
symmetric. Both are external set structures, since $|M|$ is a set.
The ordinary forcing quotient over nontransitive models is
described in \cite[Section~4, p.~12]{Williams2019}.

\begin{lemma}[Truth for the symmetric quotient]
\label{lem:arbitrary-model-symmetric-truth}
For each standard first-order formula $\varphi$ and internal
hereditarily symmetric names $\vec\tau$,
\[
 W\models\varphi([\vec\tau]_G)
 \quad\Longleftrightarrow\quad
 (\exists p\in G)\ M\models
      p\forces_{\HS_{\mathcal S}}\varphi(\vec\tau).
\]
In particular $W\models\ZF$.
\end{lemma}

\begin{proof}
Induct externally on the finite syntax of the standard formula,
using atomic formulas, conjunction, negation, and existential
quantification. The atomic case is
Equations~\eqref{eq:internal-name-equivalence} and
\eqref{eq:internal-name-membership}; atomic symmetric forcing
agrees with ordinary forcing. Conjunction uses a common
strengthening of two conditions in $G$ and monotonicity.

For negation, the set of conditions deciding the particular
formula is an internally definable dense subset of $\mathbb Q$.
It belongs to $M$ by the Separation instance for its defining
forcing predicate. Genericity gives a deciding condition in $G$.
The induction hypothesis excludes the wrong decision, and
directedness excludes conditions in $G$ forcing opposite answers.

Suppose $p\in G$ symmetrically forces an existential statement.
The set of local witness conditions in
Equation~\eqref{eq:symmetric-existential-density}, together with
the conditions incompatible with $p$, is another internally
definable dense set. Genericity meets its first part. The internal
witness is an actual element of $|M|$ which $M$ regards as an HS
name. Its equivalence class is a witness in $W$ by induction.
Conversely, a witness in $W$ has an internal HS representative;
apply induction to that representative and the elementary
existential forcing rule.

This proves the truth equivalence for every standard formula.
Lemma~\ref{lem:symmetric-forcing-schema}, applied in $M$ to each
standard $\ZF$ axiom instance, now gives $W\models\ZF$.
The proof never evaluates a name by external recursion over
$E^M$. Nor does it require a truth predicate in $M$ for all
formulas, including its possible nonstandard formula codes.
\end{proof}

The membership-closed assertion used above can also be seen
directly in this quotient. If $[\sigma]_G E_G[\tau]_G$ and
$\tau$ is internally HS, the atomic forcing rule and genericity
provide an active member name $\rho$ of $\tau$ with
$[\rho]_G=[\sigma]_G$. Heredity makes $\rho$ internally HS.
Thus $W$ and $M[G]$ have the same members of each $W$-object.
This relational form of transitivity does not make $E_G$ an
externally well-founded relation.

\subsection{Applying the construction}

The internal forcing assertion
\eqref{eq:construction-internal-forcing} can now be used in a
countable choice ground whether or not its membership relation is
externally well founded.

\begin{proof}[Proof of Theorem~\ref{thm:main-consistency}]
Assume $\Con(\ZF)$. By
Lemma~\ref{lem:constructible-consistency},
$\ZFC+V=L$ is consistent. Completeness and downward
L\"owenheim--Skolem give a countable model $M$ of this theory.
No transitivity or external well-foundedness is asserted.

Interpret the construction inside $M$ and choose an external
$M$-generic filter on its internal poset as above. Form the
symmetric quotient $W$ using
Equations~\eqref{eq:internal-name-equivalence} and
\eqref{eq:internal-name-membership}.
Lemma~\ref{lem:arbitrary-model-symmetric-truth} gives full $\ZF$
in $W$. The same lemma, applied to
Equation~\eqref{eq:construction-internal-forcing}, gives
\[
 W\models\ZF+\PP+\neg\AC+\ACwo.
\]
Soundness therefore yields
$\Con(\ZF+\PP+\ACwo+\neg\AC)$, as required.

Equivalently, one can extract the consistency implication directly
from finite proofs. A purported contradiction proof from the
target theory uses only finitely many $\ZF$ instances.
Translate its lines into symmetric forcing. The logical forcing
rules preserve its inferences;
Lemma~\ref{lem:symmetric-forcing-schema} supplies its $\ZF$
axioms; and Equation~\eqref{eq:construction-internal-forcing}
supplies $\PP$, $\ACwo$, and $\neg\AC$. This would make a nonempty
poset force a contradiction, contrary to its atomic forcing laws.
It therefore gives a contradiction proof in $\ZFC+V=L$, which
the constructible interpretation translates into one in $\ZF$.
This argument uses no stronger model-existence assumption.
\end{proof}

\subsection{The separate transitive-model conclusion}

\begin{proof}[Proof of Theorem~\ref{thm:main-transitive}]
Work in a $\ZFC$ metatheory, and let $V$ be the given
countable transitive model of $\ZFC$. Carry out
Theorem~\ref{thm:construction-summary} inside $V$, using its
parameters
\[
 \kappa=\omega_1^V,\qquad
 \mu=\bigl((2^{\aleph_0})^+\bigr)^V,\qquad
 I=(\mu^+)^V.
\]
These are the arbitrary-choice-ground parameters of the
construction.

Enumerate externally the dense subsets of the constructed poset
which belong to $V$, and recursively choose a descending
sequence meeting them. Its upward-generated filter $G$ is
$V$-generic. Here genuine well-founded recursion on ground names
is available, since $V$ is transitive. Define
\[
 W=\{\tau^G:\tau\in\HS^V\}.
\]
The set symmetric-extension theorem gives transitive models
$V\subseteq W\subseteq V[G]$ with the same ordinals, and gives
all of $\ZF$ in $W$.
Theorem~\ref{thm:construction-summary} supplies
$\PP$, $\neg\AC$, and $\ACwo$ in this same $W$, together
with its stated preservation of countable sequences from any fixed
ground set. To obtain the last assertion of
Theorem~\ref{thm:main-transitive}, let $f\in W$ be a countable
sequence of elements of $V$. Its rank in $V[G]$ is bounded by a
ground ordinal $\alpha$, since forcing preserves the ordinals.
Ranks of ground sets are absolute between these transitive models,
so the range of $f$ lies in the ground set $(V_\alpha)^V$.
The preservation assertion of Theorem~\ref{thm:construction-summary}
applied to this set gives $f\in V$. This proves the stated
transitive-model conclusion.

This is a separate model-existence deduction. The ordinary
consistency proof permits an externally ill-founded ground and
does not itself assert the existence of a transitive model.
\end{proof}
